\chapter{Technical Background and Design Choices}
\label{chap:background}

Simulation, transport, planning and visual estimation required explicit
responsibilities so that the complete system remained inspectable and testable.
The necessary ROS~2, Gazebo, rail and planning concepts lead to a comparison of
learning and control architectures and the rationale for the final layered
design.

\section{ROS 2 as the integration layer}

ROS~2 (Robot Operating System~2) is an open-source robotics software framework
that connects processes called \emph{nodes} through typed communication.
Topics carry streams, services implement request--response operations and
actions support longer tasks with feedback \citep{ros2jazzyinterfaces}. Fixed
message schemas and namespace prefixes prevent ambiguous data and name
collisions between robots; \cref{chap:multi-robot,chap:interfaces} apply these
mechanisms. Quality of Service (QoS) configures delivery properties
\citep{ros2qos}; because cameras, state and commands have different timing
needs, the implementation also checks that safety-relevant reports are fresh.

\section{Gazebo Harmonic and description formats}

Gazebo Harmonic is a simulator release \citep{gazeboharmonic}. I used it with
ROS~2 Jazzy on Ubuntu~24.04 for this project. The Simulation Description Format
(SDF) describes worlds, models, physics, plugins and sensors \citep{sdfspec}.
The Unified Robot Description Format (URDF) describes robot link--joint trees,
from which \code{robot\_state\_publisher} publishes coordinate transforms
\citep{robotstatepublisher}. \code{ros\_gz\_bridge} translates selected
messages between ROS~2 and Gazebo Transport \citep{gazeborosintegration}.
In this project, I used these tools for interface checks, teaching, repeatable
rail and safety scenarios, and labelled camera-data generation.

These purposes require different levels of fidelity. Detailed visual meshes
support cameras, and simplified collision geometry is used where contact is
modelled. The industrial arms deliberately omit collision elements, so they are
debug models rather than contact-validation models. The rail study likewise
uses repeatable kinematic poses instead of unmeasured wheel--rail forces
(\cref{chap:digital-twin,chap:transport}).

ISO~23247-1 provides an overview, terms, definitions and requirements for a
manufacturing digital-twin framework \citep{iso23247}. In this project, I used
that standard as a general framing reference and implemented a versioned
simulated environment with defined information flows, assumptions and
boundaries; visual similarity alone was insufficient.

\section{Coordinate and rail representations}

A rail path is represented by piecewise geometric segments. Each segment has
an arc-length coordinate \(s\), a physical length \(L\), and a normalised
coordinate
\begin{equation}
  \rho = \frac{s}{L}, \qquad 0 \leq \rho \leq 1 .
  \label{eq:sratio}
\end{equation}
The controller evaluates pose and heading from $s$; cubic Hermite
interpolation joins ordered computer-aided-design (CAD) samples stored as
comma-separated values (CSV) with continuous
tangents. A separate graph defines legal connections for each switch state:
geometry determines \emph{where the rail is}, while topology determines
\emph{which route is open} (\cref{fig:geometry-versus-topology}).

\begin{figure}[H]
  \centering
  \resizebox{0.96\textwidth}{!}{\begin{tikzpicture}[
  font=\sffamily\small,
  panel/.style={draw=black!25,fill=mfjalight!45,rounded corners=2mm,
    minimum width=158mm,minimum height=48mm},
  title/.style={font=\sffamily\bfseries\large,text=mfjablue},
  arr/.style={-{Latex[length=2.2mm]},thick,mfjablue},
  seg/.style={draw=mfjablue,fill=white,rounded corners=2pt,
    minimum width=22mm,minimum height=9mm,font=\ttfamily\small},
  sw/.style={diamond,draw=mfjaorange,fill=mfjaorange!12,aspect=1.7,
    minimum width=18mm,minimum height=10mm,font=\sffamily\bfseries}
]
  \node[panel] (geometry) {};
  \node[title,anchor=north west] at ([xshift=4mm,yshift=-3mm]geometry.north west)
    {Geometry: where is the shuttle?};

  \coordinate (p0) at ([xshift=16mm,yshift=2mm]geometry.west);
  \coordinate (p1) at ([xshift=-16mm,yshift=2mm]geometry.east);
  \coordinate (c0) at ($(p0)+(30mm,20mm)$);
  \coordinate (c1) at ($(p1)+(-35mm,-18mm)$);
  \draw[black!55,line width=2.1pt]
    (p0) .. controls (c0) and (c1) ..
    node[pos=0.53,coordinate] (ps) {} (p1);

  % Endpoint and local tangents follow the corresponding Bézier derivatives.
  \draw[-{Latex[length=2mm]},thick,mfjagreen]
    (p0) -- ++(18mm,12mm)
    node[right=1mm] {\(\mathbf t(0)\)};
  \draw[-{Latex[length=2mm]},thick,mfjagreen]
    (p1) -- ++(15mm,7.7mm)
    node[above left] {\(\mathbf t(1)\)};
  \draw[-{Latex[length=2mm]},thick,mfjaorange]
    (ps) -- ++(16mm,-3.2mm)
    node[pos=0.72,above=1mm] {\(\mathbf t(s)\)};
  \fill[mfjaorange] (ps) circle (1.2mm);
  \node[above=4mm of ps,font=\sffamily\small] {position \(P(s)\)};

  \fill[mfjablue] (p0) circle (1.1mm);
  \fill[mfjablue] (p1) circle (1.1mm);
  \node[below left] at (p0) {\(P(0)\)};
  \node[below right] at (p1) {\(P(1)\)};

  \coordinate (axisleft) at ([xshift=16mm,yshift=14mm]geometry.south west);
  \coordinate (axisright) at ([xshift=-16mm,yshift=14mm]geometry.south east);
  \coordinate (saxis) at (ps |- axisleft);
  \draw[mfjablue,thick] (axisleft) -- (axisright);
  \draw[mfjablue] ($(axisleft)+(0,-2mm)$) -- ($(axisleft)+(0,2mm)$);
  \draw[mfjablue] ($(saxis)+(0,-2mm)$) -- ($(saxis)+(0,2mm)$);
  \draw[mfjablue] ($(axisright)+(0,-2mm)$) -- ($(axisright)+(0,2mm)$);
  \draw[mfjaorange!80,densely dashed] (ps) -- (saxis);
  \node[below=1mm] at (axisleft) {0};
  \node[below=1mm] at (saxis) {\(s\)};
  \node[below=1mm] at (axisright) {\(L\)};
  \node[anchor=south,font=\sffamily\small,text=black!70]
    at ([yshift=1mm]geometry.south)
    {Arc length \(0\le s\le L\); normalised position \(\rho=s/L\)};

  \node[panel,minimum height=62mm,below=6mm of geometry] (topology) {};
  \node[title,anchor=north west] at ([xshift=4mm,yshift=-3mm]topology.north west)
    {Topology: which successor is legal?};

  \node[seg] (a23) at ([xshift=-47mm,yshift=8mm]topology.center) {A23};
  \node[sw,right=20mm of a23] (a3) {A3};
  \node[seg,right=30mm of a3,yshift=11mm] (a3e) {A3E};
  \node[seg,right=30mm of a3,yshift=-11mm] (a3i) {A3I};
  \draw[arr] (a23) -- node[above=1mm,font=\sffamily\small]
    {exit switch} (a3);
  \draw[arr] (a3) -- node[above,sloped,font=\sffamily\small]
    {E: exterior} (a3e);
  \draw[arr] (a3) -- node[below,sloped,font=\sffamily\small]
    {I: interior} (a3i);
  \node[anchor=south,align=center,text width=135mm,
    font=\sffamily\small,text=black!72]
    at ([yshift=5mm]topology.south)
    {After \texttt{A23}, switch \texttt{A3} selects the next public segment.\\
     Geometric proximity alone never defines a legal route.};
\end{tikzpicture}
}
  \caption[Geometry and topology as complementary rail representations.]
  {Complementary rail representations. Hermite geometry supplies
  position and tangent for pose construction; topology and switch state select
  the valid successor exposed through the public interface.}
  \label{fig:geometry-versus-topology}
\end{figure}

The ratio supports relative placement, while $s$ preserves metric distance.
Implementation and arrival semantics appear in Chapters~7--10.

\section{Symbolic task planning}

The Planning Domain Definition Language (PDDL) describes objects, predicates,
actions, preconditions and effects \citep{pddl}. PlanSys2 integrates this model
with ROS~2 \citep{plansys2}, while POPF searches for a sequence satisfying the
stated goal \citep{popf}.

A typical Room~315 planning task is to move a selected shuttle from its current
slot to a requested target slot. Symbolic planning suits this task because the
decisions are discrete and relational. The state records which shuttle occupies
the source slot, whether the target and route are free, which branch is selected,
whether devices are ready and whether the selected shuttle reached the target.
A blocker is another shuttle occupying the requested route; when
necessary, the plan first stages it temporarily on a designated interior
holding segment.

PlanSys2 supplies named actions that prepare devices and request motion. The
controller handles continuous travel; project-specific translation, supervision
and effect checks protect each step (\cref{chap:neurosymbolic}).

\section{Architectural alternatives and final design}

Vision--language--action architectures assign learning different
responsibilities. RT-2 \citep{rt2} and TinyVLA \citep{tinyvla} map observations
and language to actions; SayCan ranks predefined skills \citep{saycan};
OWL-TAMP supplies vision-language-model-generated constraints to a
task-and-motion planner \citep{owltamp}. One recent comparison found
neuro-symbolic methods outperforming VLAs on its structured long-horizon
manipulation tasks \citep{pricenotright}. I therefore treat that result as
task-specific, not as a universal ranking. \Cref{tab:architecture-comparison}
relates these patterns to Room~315.

\begin{table}[H]
\centering
\caption{Representative architectures compared by the roles assigned to
learning, planning and command execution.}
\label{tab:architecture-comparison}
\begingroup
\footnotesize
\setlength{\tabcolsep}{3pt}
\renewcommand{\arraystretch}{1.08}
\begin{tabularx}{\textwidth}{@{}
>{\raggedright\arraybackslash}p{0.15\textwidth}
>{\raggedright\arraybackslash}p{0.18\textwidth}
>{\raggedright\arraybackslash}p{0.16\textwidth}
>{\raggedright\arraybackslash}p{0.17\textwidth}
>{\raggedright\arraybackslash}X@{}}
\toprule
\textbf{Approach} & \textbf{Learned output} & \textbf{Planning role} &
\textbf{Execution path} & \textbf{Decision for Room~315}\\
\midrule
End-to-end VLA policy & Robot-action tokens or trajectories & Action selection
inside the learned policy; no separate symbolic planner & Decoded action passed
to the robot control stack & Explored but not selected; the rationale follows
the table.\\
\tablerowrule
Language/vision model with skills or a planner & Ranked skills, subgoals or planning
constraints & Skill selector or task-and-motion planner & Selected skill or
downstream controller & Related pattern; open-world constraint generation was
unnecessary here.\\
\tablerowrule
Neuro-symbolic planning with learned control & Symbolic abstractions and learned
low-level control & A PDDL
planner determines the high-level sequence & Learned controller executes each
planned PDDL action & Closest architectural family; adapted because Room~315 learns
task and scene inputs but retains an explicit controller.\\
\tablerowrule
Implemented Room~315 system & Validated task fields and paired-camera state;
neither is a plan or action & PlanSys2/POPF normally; the executive may prepend
one route-normalisation step & Software execution supervisor
publishes a typed command; the rail controller applies it &
Selected project architecture; the rationale follows the table.\\
\bottomrule
\end{tabularx}
\endgroup
\end{table}

\paragraph{Decision for Room~315.}\label{sec:ai-planning-decision}
The first exploratory architecture considered for Phase~2 used VLA-style direct
action: task language, camera frames and structured observable state would
select a compact symbolic event for software-supervisor validation and
translation, not a motor command.
An offline model was trained on planner-generated event labels, yet those traces
imitated the existing planning workflow rather than providing independent human
demonstrations or continuous robot trajectories. One learned-action run was
also excluded because target-correlated fields had leaked into its input, and
the direct-action path established no valid closed-loop result.

Room~315 instead has a finite, known decision domain: shuttle identities, slots,
directed routes, devices and admissible actions can be enumerated. The retained
V4 corpus labels scene state, the uncertain input needed by the planner. I
therefore encoded the known rules as PDDL preconditions and effects and used
PlanSys2/POPF to produce an inspectable, testable and traceable sequence.
Learning remained useful for English understanding and visual-state estimation;
neither model selects or executes an action.

Creating a leakage-free sequential corpus and independently qualifying a
learned-action policy would have added substantial data, integration and
verification work within the remaining fixed schedule. The layered design was
therefore the lower-risk fit for the available knowledge, final data and
verification method. This project-specific choice trades open-world flexibility
for manual PDDL engineering; it does not show that symbolic planning is always
superior.

Learned models interpret the task and scene. PlanSys2/POPF produces the symbolic
sequence, while deterministic software authorises commands and checks their
effects.
